% ============================================================
%  CAPÍTULO 11 – Diseño
% ============================================================
\chapter{Diseño}

\section{Arquitectura de la aplicación Flutter}

La aplicación FitLabs sigue una arquitectura orientada a pantallas (\textit{screen-based architecture}), donde cada módulo funcional corresponde a una clase \texttt{Widget} de Flutter independiente. La navegación se gestiona de forma centralizada mediante el sistema de rutas nombradas de \texttt{MaterialApp} en \texttt{main.dart}:

\begin{lstlisting}[language=Java, caption={Definición de rutas en main.dart}]
routes: {
  '/login':          (context) => const LoginScreen(),
  '/mensajes':       (context) => const MensajesScreen(),
  '/resumen':        (context) => const ResumenDiaScreen(),
  '/clientes':       (context) => const MisClientesScreen(),
  '/calendario':     (context) => const CalendarioScreen(),
  '/registrarse':    (context) => const RegistrarseScreen(),
  '/crear-rutina':   (context) => const CrearRutinaScreen(),
  '/detalle-cliente':(context) => const DetalleClienteScreen(),
  '/search-ejercicio':(context) => const SearchExerciseScreen(),
}
\end{lstlisting}

Esta estructura facilita la mantenibilidad y la incorporación de nuevas pantallas sin necesidad de modificar partes existentes del código.

\section{Sistema de diseño visual}

La identidad visual de FitLabs fue definida antes del inicio del desarrollo mediante prototipado en Figma. El diseño puede revisarse en el prototipo público disponible en \url{https://www.figma.com/design/xCVkUSQO1mjtpJQ3Z64uXH/Sin-título?node-id=0-1&t=Yd1SPcBnxRwi8qKQ-1}. Los elementos principales del sistema de diseño son:

\begin{itemize}
  \item \textbf{Paleta de colores}: Fondo degradado en tonos púrpura oscuro (\texttt{\#1A0A3C} $\to$ \texttt{\#3D2B6B}), acentos en púrpura claro (\texttt{\#9B7FD4}) y blanco para texto principal.
  \item \textbf{Tipografía}: RubikVinyl para el logotipo y títulos de alto impacto; fuente de sistema para el cuerpo del texto.
  \item \textbf{Componentes}: Tarjetas con fondo semitransparente y bordes redondeados, botones con gradiente, campos de texto con borde sutil, iconos minimalistas.
  \item \textbf{Barra de navegación inferior}: 4 ítems (Inicio, Clientes, Calendario, Mensajes) con iconos personalizados.
\end{itemize}

\section{Descripción de las pantallas principales}

\subsection{Login y Registro}

Las pantallas de autenticación presentan un diseño oscuro con el logotipo de FitLabs en la parte superior. El formulario de registro recoge nombre completo, fecha de nacimiento, email, contraseña y teléfono, con validación de campos antes de la llamada a Supabase Auth.

\begin{figure}[H]
  \centering
  \includegraphics[width=0.30\textwidth]{pantalla_login.png}
  \hspace{1cm}
  \includegraphics[width=0.30\textwidth]{pantalla_registro.png}
  \caption{Pantallas de Login (izq.) y Registro (dcha.) de FitLabs.}
  \label{fig:login_registro}
\end{figure}

\subsection{Resumen del Día}

Dashboard principal del entrenador. Muestra contadores de sesiones realizadas, pendientes y mensajes sin leer, junto con las próximas sesiones del día en tarjetas expandibles. Incluye acceso rápido a las opciones de añadir cliente, crear rutina y revisar pagos.

\begin{figure}[H]
  \centering
  \includegraphics[width=0.35\textwidth]{pantalla_resumen.png}
  \caption{Pantalla de Resumen del Día.}
  \label{fig:resumen}
\end{figure}

\subsection{Mis Clientes}

Listado paginado de clientes del entrenador con buscador en tiempo real. Cada tarjeta de cliente muestra el avatar, el nombre, el estado de la última sesión y accesos rápidos al chat y a las rutinas asignadas.

Además del filtrado por nombre, la pantalla incorpora estados visuales para mejorar la lectura rápida de la carga de trabajo del entrenador. En la implementación final se priorizó la interacción en un máximo de dos toques para acceder a cualquier cliente, reduciendo el tiempo medio de navegación en sesiones con alta rotación de usuarios.

Desde el punto de vista de arquitectura de UI, esta pantalla concentra tres patrones clave de Flutter que luego se reutilizan en otros módulos:

\begin{itemize}
  \item \textbf{Búsqueda reactiva con estado local}: actualización inmediata de la lista a medida que el usuario escribe.
  \item \textbf{Tarjetas reutilizables}: componente visual común para representar entidades de negocio (clientes, conversaciones, sesiones).
  \item \textbf{Acciones contextuales}: botones de acceso rápido para abrir chat, ver perfil o asignar una rutina sin abandonar la pantalla.
\end{itemize}

\begin{figure}[H]
  \centering
  \includegraphics[width=0.36\textwidth]{pantalla_buscar_cliente.png}
  \caption{Pantalla de búsqueda y filtrado de clientes.}
  \label{fig:buscar_cliente}
\end{figure}

\subsection{Detalle de Cliente}

Perfil completo del cliente con datos biométricos (IMC, estado de suscripción, sesiones realizadas), gráfico de progreso de peso y listado de sesiones pasadas con valoración.

Este módulo se diseñó como una vista de decisión para el entrenador. No solo presenta información histórica, sino que también facilita la toma de decisiones sobre la progresión de carga y frecuencia de sesiones. Por ello, se agruparon los datos en tres bloques:

\begin{enumerate}
  \item \textbf{Resumen de estado}: métricas rápidas para evaluar la situación actual del cliente.
  \item \textbf{Evolución}: comparación de medidas y desempeño por periodos.
  \item \textbf{Histórico operativo}: sesiones completadas con valoración para detectar adherencia.
\end{enumerate}

\begin{figure}[H]
  \centering
  \includegraphics[width=0.35\textwidth]{pantalla_detalle_cliente.png}
  \caption{Pantalla de Detalle de Cliente.}
  \label{fig:detalle_cliente}
\end{figure}

\begin{figure}[H]
  \centering
  \includegraphics[width=0.35\textwidth]{pantalla_estadisticas_cliente.png}
  \caption{Vista de estadísticas y evolución del cliente.}
  \label{fig:estadisticas_cliente}
\end{figure}

\subsection{Crear Rutina}

Formulario de creación de rutinas con selección de cliente objetivo y buscador de ejercicios integrado. Cada ejercicio añadido permite configurar series, repeticiones y peso directamente desde la tarjeta del ejercicio.

\begin{figure}[H]
  \centering
  \includegraphics[width=0.35\textwidth]{pantalla_crear_rutina.png}
  \caption{Pantalla de Creación de Rutinas.}
  \label{fig:crear_rutina}
\end{figure}

\subsection{Calendario}

Vista de calendario mensual con los días que tienen sesiones programadas marcados visualmente. Al seleccionar un día, se desplazan las sesiones del día seleccionado en la parte inferior de la pantalla.

La finalidad principal de esta pantalla es reducir conflictos de agenda. Para ello, el diseño prioriza la percepción temporal sobre el detalle, mostrando primero la ocupación del mes y solo después el detalle del día seleccionado. Este enfoque evita sobrecargar la interfaz con información de todas las sesiones simultáneamente.

\begin{figure}[H]
  \centering
  \includegraphics[width=0.36\textwidth]{pantalla_calendario.png}
  \caption{Pantalla de calendario mensual con sesiones programadas.}
  \label{fig:calendario}
\end{figure}

\subsection{Mensajes}

Listado de conversaciones activas con buscador. El contador de mensajes no leídos se muestra como insignia roja sobre el avatar del cliente.

\begin{figure}[H]
  \centering
  \includegraphics[width=0.35\textwidth]{pantalla_mensajes.png}
  \caption{Pantalla de Mensajes.}
  \label{fig:mensajes}
\end{figure}

\subsection{Perfil del entrenador}

La pantalla de perfil centraliza la configuración personal y los datos de identidad profesional del entrenador. Desde esta vista se preparó la base para futuras ampliaciones como preferencias de notificaciones, personalización visual y gestión de plan de suscripción.

\begin{figure}[H]
  \centering
  \includegraphics[width=0.34\textwidth]{pantalla_perfil_entrenador.png}
  \caption{Pantalla de perfil del entrenador.}
  \label{fig:perfil_entrenador}
\end{figure}

\subsection{Flujo de ejecución de rutinas}

Durante la validación funcional se incluyó una vista de seguimiento de sesión para representar el estado de los ejercicios ya completados y pendientes, permitiendo cerrar entrenamientos con trazabilidad. Aunque esta pantalla forma parte de la evolución del módulo, se documenta aquí por su relevancia en el diseño de interacción.

\begin{figure}[H]
  \centering
  \includegraphics[width=0.35\textwidth]{pantalla_completar_rutina.png}
  \caption{Pantalla de seguimiento/completado de rutina.}
  \label{fig:completar_rutina}
\end{figure}

\section{Criterios de usabilidad y diseño de interacción}

Para evaluar la calidad de la interfaz se aplicaron principios heurísticos sencillos, adaptados al nivel de un proyecto DAM y orientados a uso móvil real:

\begin{itemize}
  \item \textbf{Visibilidad del estado del sistema}: uso de mensajes de confirmación, indicadores de carga y estados vacíos en listados.
  \item \textbf{Consistencia visual}: misma jerarquía de tipografías, espaciados y estilo de botones en todos los módulos.
  \item \textbf{Prevención de errores}: validación de formularios antes de enviar datos a Supabase y confirmaciones en acciones sensibles.
  \item \textbf{Reconocimiento antes que recuerdo}: iconografía estable y navegación fija inferior para reducir carga cognitiva.
  \item \textbf{Eficiencia de uso}: accesos rápidos a acciones frecuentes (crear rutina, abrir chat, ver cliente) desde pantallas principales.
\end{itemize}

\section{Decisiones de diseño orientadas a mantenibilidad}

Más allá del apartado visual, se tomaron decisiones de diseño con impacto directo en la evolución del código:

\begin{enumerate}
  \item \textbf{Widgets reutilizables}: tarjetas, cabeceras y campos de formulario con estilos homogéneos para minimizar duplicidad.
  \item \textbf{Nomenclatura de rutas explícita}: rutas nominales claras y consistentes para simplificar depuración y navegación.
  \item \textbf{Separación por pantallas}: cada módulo en su archivo para reducir acoplamiento y facilitar trabajo en paralelo del equipo.
  \item \textbf{Preparación para escalado}: estructura compatible con migración futura a gestor de estado (Provider/BLoC) sin rehacer la UI.
\end{enumerate}

\section{Decisiones de diseño destacadas}

\begin{enumerate}
  \item \textbf{Diseño oscuro como estándar}: El diseño oscuro (dark mode) fue elegido como único modo de la aplicación por su asociación con entornos deportivos profesionales y por la mejor legibilidad en exteriores bajo luz solar con dispositivos actuales.

  \item \textbf{Navegación mediante barra inferior fija}: Se descartó el uso de un cajón de navegación (\textit{drawer}) en favor de una barra inferior con 4 pestañas, alineado con las directrices de Material Design 3 para apps móviles con un número reducido de secciones principales.

  \item \textbf{Separación de responsabilidades por archivo}: Cada pantalla está implementada en su propio archivo \texttt{.dart}, sin mezclar lógica de negocio y presentación en el mismo fichero. Esto facilita la colaboración entre los dos desarrolladores y reduce los conflictos en el repositorio Git.

  \item \textbf{Figma como fuente de verdad del diseño}: Se estableció como norma que ningún cambio de diseño se implementaría en código sin ser primero aprobado en el prototipo de Figma, garantizando coherencia visual entre módulos desarrollados por distintos miembros del equipo.
\end{enumerate}
